% \ifeof is an *expandable* conditional (tex.web 509), not a stomach command.
% expl3's \ior_if_eof:N is \prg_new_conditional and runs \if_eof:w (\ifeof)
% purely by expansion; its read loops -- \ior_map_variable:NNn reading
% UnicodeData.txt -- never terminate if the test cannot be answered in the
% gullet. It was implemented only in the stomach, so \ifeof inside \message or
% a fully-expandable loop produced nothing and spun. Also checks that \read
% ends each line with the \endlinechar space (tex.web 483), which the kernel's
% line parsers rely on.
\input plain
\newwrite\wout
\immediate\openout\wout=probe.dat
\immediate\write\wout{alpha}
\immediate\write\wout{beta}
\immediate\write\wout{gamma}
\immediate\closeout\wout
\newread\rin
\openin\rin=probe.dat
% \ifeof in a gullet (\message) context: must yield N, not nothing.
\message{<start=\ifeof\rin Y\else N\fi>}
\read\rin to\a
% the trailing endlinechar space: \a is `alpha ` (letters + one space)
\message{<a=\a|>}
\read\rin to\b
\message{<b=\b|>}
\read\rin to\c
\message{<eofnow=\ifeof\rin Y\else N\fi>}
\message{<c=\c|>}
\read\rin to\d
\message{<afterlast=\ifeof\rin Y\else N\fi>}
\closein\rin
% a recursive read-until-eof driven entirely by the expandable \ifeof, the
% shape \ior_map_variable uses.
\openin\rin=probe.dat
\count0=0
\def\rdall{\ifeof\rin \else \read\rin to\line \advance\count0 by1 \expandafter\rdall\fi}
\rdall
\message{<looptotal=\the\count0>}
\closein\rin
% \ifeof on an unopened stream reads as true
\message{<unopened=\ifeof9 Y\else N\fi>}
% \read strips trailing spaces from the line (tex.web 31) before the endline
% space, so `trail   ' reads as `trail ' -- one space, not four.
\immediate\openout\wout=trail.dat
\immediate\write\wout{trail   }
\immediate\closeout\wout
\openin\rin=trail.dat
\read\rin to\t
\message{<trail=\meaning\t>}
\closein\rin
% \readline reads a line verbatim: every character catcode 12 (other) except
% a space (catcode 10). expl3's \ior_str_get:NN reads CaseFolding.txt this way.
\immediate\openout\wout=str.dat
\immediate\write\wout{a b~c}
\immediate\closeout\wout
\openin\rin=str.dat
\endlinechar=-1
\readline\rin to\s
% every non-space char is `other'; \ifcat against a known other-char is true,
% and against a letter is false -- proving the catcodes are forced.
\message{<rl=\meaning\s>}
\message{<rlother=\expandafter\ifcat\string a12\noexpand\else L\fi>}
\closein\rin
\bye
